\section{Discussion}
\label{sec:discussion}

\para{Fertility is a moderate correlate, not the dominant driver.}
The hypothesis treats fertility as the main, tokenizer-specific cause of transfer loss.
Our data describe it as a moderate correlate that is largely redundant with better predictors.
Which predictor matters most depends on the task: syntactic distance for NER, pretraining volume for topic classification, and both for NLI.
This agrees with the typology-and-size account of \citet{lauscher2020zero} and with the findings of \citet{wu2020all} on low-resource languages.

\para{Fertility is not irrelevant either.}
The within-language contrast and the intervention both point to a real effect of fragmentation.
The reconciliation we find most plausible is that fertility mixes two things: fragmentation that the model was never trained to handle, which is harmful, and fragmentation that reflects ordinary morphology in a well-resourced language, which is mostly harmless.
A single ratio cannot tell them apart, which would explain why Finnish and Malayalam can both be heavily fragmented and have very different outcomes.
This is a hypothesis generated by our results, not something they establish.

\para{Surprises.}
First, \xlmr never reaches the hypothesized danger zone for any whitespace-segmented language, so a threshold defined on fertility is relative to a tokenizer and does not carry over to another.
Second, the choice of fertility definition changes the \sib answer from weak to moderately strong.
Third, over-segmented English does far worse than natural languages at the same fertility.

\para{Practical implications.}
Fertility alone is a weak screening signal for transfer risk.
Syntactic distance and pretraining volume are more informative, and for NER so is script.
Practitioners should not use the 2.5$\times$ threshold as a rule.

\para{Limitations.}
\begin{itemize}[leftmargin=*,itemsep=0pt,topsep=0pt]
    \item \emph{Two models.} Both are base-size encoders. Our results do not speak to decoder-only language models or to byte- and character-level models.
    \item \emph{Observational design.} Fertility, size, typology and script are correlated across languages, and with 67 languages (12 for \xnli) partial effects have wide intervals. The \xnli cells cannot rule out a fertility \rsq of 0.4--0.6 (upper bounds 0.60 and 0.40), although the point estimates are 0.16 and 0.00.
    \item \emph{Proxies.} \mbert's pretraining size uses 2023 Wikipedia dump sizes, not the 2018 dump the model was trained on, and \xlmr's uses the public CC-100 re-creation. We computed \uriel distances ourselves from kNN-imputed features. We compiled \xlmr's language list from our recollection of the original paper and excluded Haitian and Yoruba as uncertain.
    \item \emph{Reduced \xnli training.} Training on 100k examples lowers absolute accuracy; we did not check whether it changes relative drops.
    \item \emph{Noisy cells.} \sib has 204 test sentences per language and ten \wikiann languages have 100. \wikiann labels are automatically derived and uneven in quality, which adds language-specific noise unrelated to any predictor.
    \item \emph{Whitespace fertility.} It is undefined for four languages and is a crude proxy for words in others; the sensitivity to the fertility definition in \secref{sec:results} shows that this matters.
    \item \emph{Intervention validity.} The intervention creates out-of-distribution segmentations, so it measures the cost of unfamiliar fragmentation. It does not reproduce what a naturally high-fertility language experiences.
    \item \emph{One recipe.} We use one learning rate per task and no per-language tuning. Seed variance is small relative to between-language variance, but we did not explore recipe sensitivity.
    \item \emph{Literature coverage.} We did not conduct a systematic literature review, so our positioning relative to prior work may be incomplete.
\end{itemize}
